\section{Introduction}
\label{sec:introduction}

A bank that uses an internal model to set capital, provisions or limits must show a second line of defence that the model works, and the people who do that work, the model validators, need three things at once: a catalogue of tests that supervisors recognise, a set of reference calculations against which a bank's own numbers can be compared, and a way of knowing that the tests and the references are themselves correct. The European Central Bank fixes the first for the internal ratings based credit models of the banks it supervises: its instructions for reporting validation results prescribe, for the probability of default, a Jeffreys back-test per rating grade, a comparison of the area under the ROC curve with its value at development, the matrix weighted bandwidth and a monotonicity test of the migration matrix, and a concentration test on the grade distribution; for the loss given default and the credit conversion factor they prescribe paired $t$-tests and a generalised area statistic~\citep{ecb2019}. The Basel Committee fixes the second for the interest rate risk in the banking book: six shocks of the yield curve with sizes by currency and a standardised economic value of equity, with a worked example in the standard~\citep{bcbssrp31,bcbsd368}. The reference calculation of the capital of a loan portfolio is a family of models that dates from the late 1990s, from the one-factor structural model and its large-portfolio limit to the actuarial CreditRisk+ model and its recursions~\citep{gordy2003,gordy2000,haaf2004}.

The statistics are public and short, and the difficulty is the third need. A validator who writes the Jeffreys test as a call to a Beta distribution function has written one line, and the line is correct only if the prior is the one the instructions name, if the portfolio test uses the obligor-weighted mean of the grade probabilities and not their unweighted mean, and if the oriented tail is the one that signals an underestimated probability. A portfolio engine has far more places to be wrong: a seeded sampler that draws in chunks, a recursion on a grid of loss units, an importance sampler whose weights must average to one. Each has a failure that no single test reveals, and the usual defence, a unit test written by the same person against a number the same person computed, shares the author's misreading with the code.

We describe and measure an answer to the third need as it is built into one open package, the engine \texttt{riskvalidation}, release \inventoryVersion{}, which holds \testsTotal{} registered tests in \familiesTotal{} families, each returning one result record with its statistic, null hypothesis, primary source and policy; five of the ECB statistics; a copula and CreditRisk+ portfolio engine with Euler contributions, a granularity adjustment and an importance sampler; and an IRRBB engine with the six shocks, the standardised economic value of equity and the Chilean currency rule. The proposed method is a protocol for holding such an engine to references that do not share its code: analytic anchors, published test vectors, an independent transcription of the specification, a seeded regeneration of a committed result, a simulator with a known truth, and an independent implementation by another library; it is measured by planting defects in the engine and recording which layer of the protocol finds each.

The contributions are the following.
\begin{itemize}
\item A map of which of the \ecbSections{} statistics that the ECB instructions fix the engine implements as specified (\ecbImpl{}), which it implements in a related but different form (\ecbRelated{}) and which it does not implement (\ecbAbsent{}), obtained from the engine's registry and not from its documentation (Section~\ref{sec:results}).
\item Differential tests of the five implemented statistics against an independent transcription of the instructions on \ecbInputs{} random inputs each, and a comparison of the Jeffreys test with the nearest generic tests on the \cCells{} grades of eight credit-card models of the committed holdout of a real portfolio.
\item A reconciliation of the portfolio engine to four analytic anchors, an independent simulator and an independent recursion, to its own committed one-million-scenario rung and to a simulator with a known truth, and of the IRRBB engine to the printed worked example and to QuantLib on \irrLegs{} committed leg and shock pairs.
\item A measurement, by \ablMutants{} planted defects, of what each verification layer contributes, including the finding that a check built to detect regressions flags changes that are not defects, and a power calculation that says how large a defect an independent simulation can see.
\item One result that does not hold: the maximum-likelihood interval of the asset correlation falls outside the binomial band of its own simulator in \gridMlFail{} of \gridMlCells{} cells.
\end{itemize}

The main result is that every implemented component agrees with its independent reference to the accuracy of the reference, that the one hypothesis about estimator coverage fails in a cell, and that the five layers are not interchangeable: \ablAny{} of the \ablReal{} defects that change a result are found by at least one layer and \ablOnly{} by exactly one. The absence claim that motivated the package, that no open toolkit implements these tests as specified, is not tested here beyond a dated survey of \pypiFound{} toolkits and a scan of \distScanned{} installed distributions, and it is stated with that limit (Section~\ref{sec:limitations}). The paper is organised as follows: Section~\ref{sec:related-work} places the work among the published tests, engines and toolkits; Section~\ref{sec:problem} states the problem; Section~\ref{sec:proposed-method} gives the engine's statistics and the protocol; Section~\ref{sec:theory} proves what can be proved; Sections~\ref{sec:experimental-design} to~\ref{sec:ablation} give the design, the results and the ablation; and Sections~\ref{sec:discussion} to~\ref{sec:conclusions} close.
